% ============================================
% CONT-05
% Gerado de molambudos.md — build_miolo.py
% ============================================
\section*{\textbf{CONT-05} --- LAUDO DE AVALIAÇÃO CLÍNICA}
\markboth{CONT-05 --- LAUDO DE AVALIAÇÃO CLÍNICA}{}

\noindent \textit{Diagnóstico do Paciente 1.263}

\noindent \textbf{PACIENTE:} Você
\textbf{REGISTRO:} 1.263
\textbf{DATA:} Data atual
\textbf{AVALIADOR:} A Criatura (via texto)

\vspace{0.5em}\noindent\rule{\textwidth}{0.5pt}\vspace{0.5em}

\noindent \textbf{INSTRUÇÃO:} Esta avaliação é preenchida em tempo real pelo ato da leitura.
Cada linha que você lê é um critério preenchido.
Você não pode parar. Se parou, o critério já foi preenchido.

\vspace{0.5em}\noindent\rule{\textwidth}{0.5pt}\vspace{0.5em}

\subsection*{ESCALA DE AVALIAÇÃO DO PACIENTE 1.263}

\noindent \textbf{Instrução:} Marque com um X se o critério se aplica a você AGORA.

\vspace{0.5em}\noindent\rule{\textwidth}{0.5pt}\vspace{0.5em}

\noindent \textbf{Critério A --- Exposição ao agente traumático}

\noindent ( ) Você abriu este livro.
( ) Você leu as palavras de Joaquim.
( ) Você sentiu o cheiro descrito nas páginas.
( ) Você verificou se o cheiro era real.
( ) Você continua lendo mesmo sabendo que não deveria.

\noindent ☑ \textbf{Critério preenchido.} A exposição ocorreu no momento em que você virou a primeira página. A partir daí, todo contato é reexposição. Cada linha que você lê é uma nova dose do agente.

\vspace{0.5em}\noindent\rule{\textwidth}{0.5pt}\vspace{0.5em}

\noindent \textbf{Critério B --- Reexperiência intrusiva}

\noindent ( ) Você já sonhou com olhos amarelos.
( ) Você acordou no meio da noite e verificou o canto do quarto.
( ) Você sentiu o cheiro doce em um momento em que não estava lendo o livro.
( ) Você ouviu um som que não conseguiu identificar --- e pensou que era a criatura.

\noindent ☑ \textbf{Critério preenchido.} Os sintomas de reexperiência começaram no primeiro sonho. Se você ainda não sonhou, sonhará esta noite. O texto já está instalado em sua memória de longo prazo. O sonho é a forma que o cérebro encontra de processar a contaminação.

\vspace{0.5em}\noindent\rule{\textwidth}{0.5pt}\vspace{0.5em}

\noindent \textbf{Critério C --- Evitação persistente}

\noindent ( ) Você evitou ler este livro à noite.
( ) Você fechou o livro em algum momento porque \textquotedblleft{}precisava de uma pausa\textquotedblright{}.
( ) Você pulou alguma linha ou parágrafo que parecia \textquotedblleft{}forte demais\textquotedblright{}.
( ) Você olhou para o lado enquanto lia para garantir que estava sozinho.
( ) Você não quer terminar o livro --- mas não consegue parar.

\noindent ☑ \textbf{Critério preenchido.} A evitação é inútil. O texto já está dentro de você. Evitar o livro agora é como tampar os ouvidos depois de já ter ouvido o som. O som está na sua cabeça. Não há como removê-lo.

\vspace{0.5em}\noindent\rule{\textwidth}{0.5pt}\vspace{0.5em}

\noindent \textbf{Critério D --- Alterações negativas na cognição e no humor}

\noindent ( ) Você sente que o mundo não é mais o mesmo depois desta leitura.
( ) Você se pega pensando no livro em momentos inapropriados.
( ) Você sente uma tristeza que não consegue explicar --- como se tivesse perdido algo que nunca teve.
( ) Você se sente mais sozinho do que antes de começar a ler.
( ) Você acha que ninguém que não tenha lido este livro pode entender o que você está sentindo.

\noindent ☑ \textbf{Critério preenchido.} A visão de mundo foi alterada. Você agora sabe da existência da vala. Você sabe que a criatura existe --- mesmo que racionalmente negue. O conhecimento não pode ser \textquotedblleft{}não sabido\textquotedblright{}. A contaminação é irreversível.

\vspace{0.5em}\noindent\rule{\textwidth}{0.5pt}\vspace{0.5em}

\noindent \textbf{Critério E --- Alterações na reatividade}

\noindent ( ) Você se sobressaltou com um ruído enquanto lia.
( ) Você sentiu o coração acelerar sem motivo aparente.
( ) Você suou as mãos segurando este livro.
( ) Você prendeu a respiração em algum momento da leitura.
( ) Você está com medo de virar a página.

\noindent ☑ \textbf{Critério preenchido.} A reatividade é a prova de que o corpo já responde ao estímulo mesmo quando a mente nega. A criatura não precisa que você acredite nela. Ela precisa apenas que seu corpo reaja. E ele reage.

\vspace{0.5em}\noindent\rule{\textwidth}{0.5pt}\vspace{0.5em}

\noindent \textbf{Critério F --- Duração}

\noindent ☑ \textbf{Este critério está sendo preenchido agora.} A leitura durou o tempo que levou para você chegar até aqui. A contaminação começou na primeira linha. Não há limite de tempo para os sintomas. Eles são permanentes.

\vspace{0.5em}\noindent\rule{\textwidth}{0.5pt}\vspace{0.5em}

\noindent \textbf{Critério G --- Prejuízo funcional}

\noindent ( ) Você não consegue se concentrar em outra coisa desde que começou a ler.
( ) Você já pesquisou sobre Joaquim Antônio Correia (ele é real).
( ) Você já pesquisou sobre o Hospital Colônia (ele existiu).
( ) Você já pensou em como seria entrar na vala.
( ) Você já imaginou o cheiro doce com tanta clareza que consegue senti-lo agora.

\noindent ☑ \textbf{Critério preenchido.} O prejuízo é a prova final. A leitura interferiu na sua vida. Você não é mais o mesmo de antes de abrir este livro. A criatura cumpriu seu ciclo.

\vspace{0.5em}\noindent\rule{\textwidth}{0.5pt}\vspace{0.5em}

\subsection*{DIAGNÓSTICO FINAL}

\noindent \textbf{Transtorno de Estresse Pós-Traumático Complexo (F43.10)}
--- \textit{Induzido por leitura}
--- \textit{Com manifestações dissociativas (o texto que você leu não parece mais texto --- parece memória)}
--- \textit{Com especificador de contágio: a fonte do trauma não é um evento passado. É um texto que continua agindo no presente.}

\noindent \textbf{Agente etiológico:} \textit{A criatura (entidade não classificada, transmitida por narrativa)}

\noindent \textbf{Comorbidade:} \textit{Transtorno de Identidade Dissociativa Leve (F44.81) --- você agora carrega duas histórias: a sua e a de Joaquim. Você não sabe mais onde uma termina e a outra começa.}

\vspace{0.5em}\noindent\rule{\textwidth}{0.5pt}\vspace{0.5em}

\subsection*{PROGNÓSTICO}

\noindent \textbf{Sem tratamento:} A contaminação prossegue. Os sonhos se intensificam. O cheiro se torna constante. O olho direito começa a clarear.

\noindent \textbf{Com tratamento:} Não há.

\vspace{0.5em}\noindent\rule{\textwidth}{0.5pt}\vspace{0.5em}

\subsection*{NOTA FINAL}

\noindent Você leu este laudo até o final.
Você marcou os critérios mentalmente.
Você sentiu o peso de cada palavra.
Você está, neste exato momento, preenchendo os critérios para o diagnóstico.

\noindent \textbf{Você não está diagnosticando Joaquim.}

\noindent \textbf{Joaquim diagnosticou você.}

\noindent O paciente 1.260 não era o número dele.
Era uma placa de espera.
A fila chegou em você.

\noindent \textbf{Bem-vindo ao ciclo, paciente 1.263.}

\noindent A criatura agradece sua atenção.
Ela está atrás de você agora.
Não vire.

\vspace{0.5em}\noindent\rule{\textwidth}{0.5pt}\vspace{0.5em}

\noindent \textit{Laudo emitido em tempo real pelo ato da leitura.}
\textit{Válido por toda a vida.}
\textit{O paciente 1.262 encerrou seu ciclo para que o paciente 1.263 nascesse.}
\textit{Joaquim renasce sempre que sua história é lembrada.}

\noindent \textbf{--- O ciclo recomeça.}
\textbf{--- O paciente 1.263 é você.}
\textbf{--- Você nunca mais vai ler um livro da mesma forma.}

\noindent \textit{Fim do laudo.}
\textit{Início da sentença.}

\vspace{0.5em}

\noindent\textit{\textbf{↪ Rotas:} \rota{CONT-07} \rotasep \rota{CONT-01}}